\chapter{Statistiques : position et lecture critique}
\section*{Exercice 18.1 -- Indicateurs}
Pour $1,2,2,3,5,5,7,8,9,10$, calculer moyenne, médiane, quartiles, étendue et écart interquartile avec les conventions du cours.

\section*{Exercice 18.2 -- Moyenne de groupes}
Un groupe de $12$ élèves a une moyenne de $14$, un groupe de $18$ a une moyenne de $11$. Donner la moyenne globale.

\section*{Exercice 18.3 -- Transformation}
Chaque mesure d'une série de moyenne $7$ est transformée par $y=3x-2$. Quelle est la nouvelle moyenne ? Justifier.

\section*{Exercice 18.4 -- Ajout extrême}
Ajouter $100$ à la série de l'exercice 18.1. Recalculer moyenne et médiane, puis commenter.

\section*{Exercice 18.5 -- Lire un diagramme}
Une catégorie compte $45$ personnes sur $180$. Quel angle occupe-t-elle dans un diagramme circulaire, puis semi-circulaire ?

\section*{Exercice 18.6 -- Mêmes centres ?}
Comparer les séries $A=(4,5,6)$ et $B=(0,5,10)$. Peut-on les dire identiques parce qu'elles ont même moyenne et même médiane ?
